\medskip

\emph{Cet exercice est un questionnaire à choix multiples.\\
Pour chacune des questions suivantes, une seule des quatre réponses proposées est exacte. Une réponse exacte rapporte un point. Une réponse fausse, une réponse multiple ou l'absence de réponse à une question ne rapporte ni n'enlève de point. Pour répondre, indiquer sur la copie le numéro de la question et la lettre de la réponse choisie. Aucune justification n'est demandée.}

\bigskip

L'espace est rapporté à un repère orthonormé \Oijk.

\medskip
On considère :

les points A$(-1~;~-2~;~3)$, B$(1~;~-2~;~7)$ et C$(1~;~0~;~2)$;

la droite $\Delta$ de représentation paramétrique : \quad $\left\{\begin{array}{l}x=1-t \\ y=2 \\ z=-4+3 t\end{array}\right.$, où $t \in \R$;

le plan $\mathcal{P}$ d'équation cartésienne :\quad $3 x+2 y+z-4=0$;

le plan $\mathcal{Q}$ d'équation cartésienne : \quad $-6 x-4 y-2 z+7=0$.

\begin{enumerate}[itemsep=5mm]
	\item Lequel des points suivants appartient au plan $\mathcal{P}$ ?

	\begin{tabularx}{\linewidth}{*{4}{>{\bfseries}p{5mm}@{}X}}
		a.& $R(1~;~-3~;~1)$;&
		b.& $S(1~;~2~;~-1)$;&
		c.& $T(1~;~0~;~1)$;&
		d.& $U(2~;~-1~;~1)$.\\
	\end{tabularx}

	\item Le triangle ABC est :

	\begin{tabularx}{\linewidth}{*{2}{>{\bfseries}p{5mm}@{}X}}
		a.& équilatéral ;&
		b.& rectangle isocèle ;\\
		c.& isocèle non rectangle ;&
		d.& rectangle non isocèle.\\
	\end{tabularx}

	\item La droite $\Delta$ est :

	\begin{tabularx}{\linewidth}{*{2}{>{\bfseries}p{5mm}@{}X}}
		a.& orthogonale au plan $\mathcal{P}$ ;&
		b.& sécante au plan $\mathcal{P}$ ;\\
		c.& incluse dans le plan $\mathcal{P}$ ;&
		d.& strictement parallèle au plan $\mathcal{P}$.\\
	\end{tabularx}

	\item On donne le produit scalaire $\vect{\mathrm{BA}} \cdot \vect{\mathrm{BC}}=20$.

	Une mesure au degré près de l'angle $\widehat{\mathrm{ABC}}$ est :

	\begin{tabularx}{\linewidth}{*{4}{>{\bfseries}p{5mm}@{}X}}
		a.& $34^{\circ}$;&
		b.& $120^{\circ}$;&
		c.& $90^{\circ}$;&
		d.& $0^{\circ}$.\\
	\end{tabularx}

	\item L'intersection des plans $\mathcal{P}$ et $\mathcal{Q}$ est :

	\begin{tabularx}{\linewidth}{*{2}{>{\bfseries}p{5mm}@{}X}}
		a.& un plan ;&
		b.& l'ensemble vide ;   \\
		c.& une droite ;&
		d.& réduite à un point.\\
	\end{tabularx}
\end{enumerate}
